\section{Related Work}

In accordance with an interdisciplinary research methodology balancing approximately 50\% Artificial Intelligence technology and 50\% environmental health domain know-how, we categorize prior studies from foundational statistical and regulatory methods to advanced machine learning and tool-augmented agent frameworks.

\subsection{Methods: AI Technology and Environmental Domain Know-How (M)}
\subsubsection{AI Technology---From Naive Baselines to Tree Ensembles, XAI, and ReAct Agents}
In time-series forecasting, classical baselines such as Persistence ($\hat{y}_{t+h} = y_t$) and Seasonal Naive ($\hat{y}_{t+h} = y_{t+h-24}$) serve as indispensable reference points. Large-scale empirical evidence from the M5 forecasting competition \cite{makridakis2022} and Hewamalage et al.'s forecast evaluation taxonomy \cite{hewamalage2023} demonstrate that complex ML models must be evaluated via rolling-origin splits and skill scores relative to naive baselines to avoid illusory gains. In air quality forecasting specifically, Chianese and Riccio \cite{chianese2026} synthesized 409 studies and concluded that no single architecture universally dominates across horizons and monitoring regimes; rather, leakage-aware validation \cite{lu2026} and feature availability at the forecast origin govern deployment reliability. Joseph and Ganapathy \cite{joseph2025} further observed that L2-regularized linear models (Ridge) and tree ensembles (Random Forest, XGBoost) frequently match deep recurrent or transformer models on tabular station histories while providing superior computational efficiency. For model interpretability, Lundberg and Lee's SHAP framework \cite{lundberg2017} is widely applied to tree models, yet Houdou et al. \cite{houdou2024} emphasized in a systematic review that SHAP quantifies local/global model feature attribution rather than physical atmospheric causality. At the interaction layer, Yao et al. \cite{yao2023} introduced ReAct to interleave reasoning traces with external tool calls, Gao et al. \cite{gao2025} demonstrated Instructor--Worker LLM coordination for wildfire air quality analysis, and Nie and Liu \cite{nie2025} alongside OWASP LLM01 (2026) \cite{owasp2026} established that environmental LLMs must operate as rule-constrained assistants with strict tool allowlists to prevent prompt injection and excessive agency.

\subsubsection{Domain Know-How---Exposure Hierarchy, Regulatory AQI, and Decision Utility}
At the public health foundation, the WHO Global Air Quality Guidelines \cite{who2021,who2024} define short- and long-term reference thresholds for criteria pollutants while noting that population guidelines are not individual clinical diagnostic criteria. In exposure science, Steinle et al. \cite{steinle2018} formalized the distinction among: 1) fixed-station ambient concentration $C(t)$ ($\mu\text{g/m}^3$); 2) time-integrated exposure proxy $\sum_i C(t_i)\Delta t_i$ ($\mu\text{g}\cdot\text{h/m}^3$); and 3) true personal inhaled dose, which requires microenvironmental infiltration, street-canyon geometry, and minute ventilation modeling. Both U.S. EPA sensor guidance \cite{epa2024} and Lo et al.'s Taiwanese spatiotemporal $\text{PM}_{2.5}$ study \cite{lo2025} confirm that sparse regulatory stations cannot validate street-level clean-route navigation without independent mobile or satellite-LUR calibration. Under Taiwan's regulatory standard, MOENV's 2025 revised AQI \cite{moenv2025aqi} aggregates six pollutants over pollutant-specific moving averages into six tiers with differentiated guidance for general and sensitive populations. Crucially, Berlinghieri et al. \cite{berlinghieri2026} evaluated six hourly $\text{PM}_{2.5}$ forecast systems for daily outdoor planning (``whether to go out'' vs. ``when to go out''), proving that lower point-forecast MAE does not monotonically yield superior activity window selection.

\begin{table*}[!t]
\centering
\footnotesize
\setlength{\tabcolsep}{4pt}
\caption{Organized Comparison and SWOT Judgement of Representative Related Work (AI Technology $\times$ Domain Know-How)}
\label{tab:literature}
\begin{tabular}{L{1.4cm}L{2.4cm}L{2.6cm}L{2.5cm}L{4.0cm}L{3.7cm}}
\toprule
\textbf{Ref. / Yr} & \textbf{Problem / Scope} & \textbf{AI / Domain Method} & \textbf{Dataset / Evidence} & \textbf{Key Findings \& Strengths (S/O)} & \textbf{Limitations \& Threats (W/T)} \\
\midrule
\cite{chianese2026} 2026 & Cross-regime air quality forecasting & Systematic mapping \& validation taxonomy & 409 air-quality forecasting papers & No universal best model; validation design and leakage prevention dominate performance & Literature synthesis; not an empirical deployment on Taiwan MOENV stations \\
\cite{hewamalage2023} 2023, \cite{lu2026} 2026 & Time-series forecast evaluation pitfalls & Rolling-origin split, leakage control, skill scores & Forecasting methodology \& urban $\text{PM}_{2.5}$ & Chronological splitting and naive baseline comparisons prevent inflated accuracy claims & Methodological guidelines; require integration with regulatory decision rules \\
\cite{makridakis2022} 2022 & Complex ML vs. naive baselines & M5 competition tree ensembles vs. naive & 42,840 hierarchical retail time series & Gradient boosting excels when rigorously benchmarked against seasonal/persistence baselines & Retail sales domain; cannot be extrapolated directly to $\text{PM}_{2.5}$ pollution spikes \\
\cite{lundberg2017} 2017, \cite{houdou2024} 2024 & Explainable ML in air pollution & SHAP additive feature attribution review & Tree and linear air quality models & Improves transparency of local and global feature contributions for non-experts & Quantifies model behavior only; must not be misinterpreted as physical causality \\
\cite{yao2023} 2023, \cite{owasp2026} 2026 & Tool-augmented LLM reasoning \& safety & ReAct prompting \& OWASP LLM01 defense & Interactive QA \& GenAI security standards & Interleaving reasoning with allowlisted tools curbs hallucination and prompt injection & Unconstrained agency or untrusted API payloads can still hijack agent outputs \\
\cite{nie2025} 2025, \cite{gao2025} 2025 & LLMs in environmental decision workflows & Supervised assistant \& Instructor--Worker LLM & Environmental modeling \& 2025 LA wildfires & LLMs effectively orchestrate environmental queries when supervised by domain models & Case-specific or conceptual; lacks Taiwan MOENV station-level backtesting \\
\cite{who2021} 2021, \cite{steinle2018} 2018 & Public health burden \& exposure dynamics & WHO AQG \& time-activity exposure hierarchy & Global epidemiology \& microenvironment data & Distinguishes ambient concentration ($\mu\text{g/m}^3$), exposure proxy ($\mu\text{g}\cdot\text{h/m}^3$), and inhaled dose & Fixed station readings cannot represent individual inhaled dose or street gradients \\
\cite{berlinghieri2026} 2026 & Hourly $\text{PM}_{2.5}$ forecasts for daily planning & Point error vs. activity scheduling utility & Six operational forecast systems (U.S. wildfires) & Demonstrates that lower overall MAE does not automatically optimize ``when to go out'' & Evaluated on U.S. wildfire seasons; local Taichung station skill requires verification \\
\cite{moenv2025aqi} 2025, \cite{lo2025} 2025 & Taiwan AQI rules \& $\text{PM}_{2.5}$ spatiotemporal ML & MOENV 2025 sub-index rules \& LUR+XGBoost & Taiwan MOENV datasets (432, 13, 488, 07) & Establishes official multi-pollutant AQI tiers and local Taiwanese $\text{PM}_{2.5}$ characteristics & Monthly LUR or static AQI lookup does not provide 1--6h auditable agent comparisons \\
\textbf{Ours (v3)} & \textbf{Station-scale 1/3/6h $\text{PM}_{2.5}$ \& activity agent} & \textbf{ReAct allowlist + ML baselines + 2025 AQI rules} & \textbf{Taichung 6 MOENV stations (123,055 valid h)} & \textbf{Cryptographic W1 data audit, dual-track AQI/$\text{PM}_{2.5}$ separation, fail-closed abstention} & \textbf{Scoped strictly to station scale; W2--W6 ML/Agent scores in active evaluation} \\
\bottomrule
\end{tabular}
\end{table*}

\subsection{New Proposed Trend, Deep Research Synthesis, and Organized Comparison (N, O)}
Building on our Deep Research investigation, AirGuard AI is anchored by five core domain mindsets (Q1): 1) \emph{Exposure context over static visualization}; 2) \emph{Spatiotemporal fidelity and three-timestamp provenance}; 3) \emph{Baseline-first and leakage-safe validation}; 4) \emph{Separation of point forecast accuracy and decision utility}; and 5) \emph{Human-centered, fail-closed agent guardrails}. Furthermore, we address three major professional controversies (Q2):
\begin{enumerate}
    \item \emph{Complex Deep Learning vs. Simple Baselines:} Rather than assuming deep neural networks outperform simpler models on single-station series, we require Ridge, Random Forest, and XGBoost to prove positive skill against Persistence and 24-hour Seasonal Naive on identical valid sample pairs.
    \item \emph{Average Forecast Error vs. Actionable Activity Decisions:} Because adjacent candidate windows often differ by less than model validation error \cite{berlinghieri2026}, we pair point-error metrics (Track A) with uncertainty-overlap abstention and decision-regret evaluation (Track B).
    \item \emph{Autonomous Agent vs. Rule-Constrained Pipeline:} To satisfy ten deep domain reasoning checks (Q3: Why single-hour $\text{PM}_{2.5} \neq$ official AQI; What-if API lags by 3 hours; Trade-off of marginal MAE gains; Failure cases of random splitting; Evidence of rule compliance; Horizon stratification; Architectural separation; Spatial station limits; Uncertainty crossing thresholds; and Non-causality of SHAP), the LLM is restricted to an eight-tool allowlist with deterministic output verification.
\end{enumerate}

Table~\ref{tab:literature} organizes the representative literature and SWOT judgement (`O`), establishing a clear four-stage progression: \emph{Existing Research} provides isolated open datasets, ML forecasters, and generic tool agents $\rightarrow$ \emph{Limitations} arise from timestamp conflation, temporal leakage, AQI/$\text{PM}_{2.5}$ confusion, and unconstrained LLM hallucination $\rightarrow$ \emph{Research Gap} lies in the absence of an end-to-end, cryptographically reproducible Taiwan MOENV station-scale framework uniting leakage-safe 1/3/6-hour baselines, versioned 2025 AQI rules, and fail-closed agent auditing $\rightarrow$ \emph{Our Project (AirGuard AI)} bridges this gap.
